% Policy Memorandum and Concept Paper. Synthesised, not narrated. File names
% live in the source table; the prose cites works by bibliography key.

\begin{flushleft}
\textbf{MEMORANDUM}\\[3pt]
\textbf{To:}\hspace{0.5cm} Committee on Energy and Commerce, U.S. House of Representatives
(Subcommittee on Health), and Members with jurisdiction over the Food and Drug
Administration\\
\textbf{From:}\hspace{0.0cm} Office of the Sponsor, H.R. 9510 (CEO Kevin Kawchak,
ChemicalQDevice)\\
\textbf{Re:}\hspace{0.45cm} VVUQ Physical AI Oncology Trial Bill: concept paper and
statement of legislative intent\\
\textbf{Date:}\hspace{0.1cm} May 30, 2026
\end{flushleft}

This memorandum accompanies and introduces H.R. 9510. It is offered in the form
that precedes a serious bill: a concept paper that fixes the idea, a statement of
why the law is needed now, and an account of the new method by which the
supporting record was assembled. The sources drawn on throughout are listed once,
with their repository locations, in Table~\ref{tab:policy-sources}; the body that
follows cites those works by name and does not repeat their paths.

\subsection{Concept and Legislative Intent}

The concept of this bill is narrow, deliberate, and new. It elevates the
verification of robot-patient interaction code above the code itself, and it
requires that verification to clear before the code is generated or executed in a
Physical AI oncology clinical trial. The instrument does not regulate a drug, a
device, or a diagnosis in the first instance; it regulates the order of
operations by which autonomous software is allowed to touch a patient. That a
statute should govern sequence rather than substance is unusual, and the reason
is specific to this moment in oncology research.

Three fronts that matured separately are now converging in the same procedure
room. In silico simulation has become accurate enough to stand in for early
clinical testing; autonomous code generation has become capable enough to write
the control software that drives a procedure; and robotic platforms have become
dexterous enough to perform that procedure on a body. The author's first VVUQ
study \cite{kawchak2026vvuq01} examined what happens when these fronts meet and
reached a conclusion that the present bill codifies: when software is generated
autonomously and then executed against a patient, the act of generation is cheap
and fast, while the act of deciding whether the generated software may run is the
work that carries the consequence. An error caught at a verification gate before
delivery costs a regenerated artifact and a few milliseconds; the same error
discovered after delivery costs a patient. That asymmetry is the entire economic
and ethical case for verification before generation, and it grows sharper as
embodiment concentrates every failure mode into one moving body.

The vocabulary this bill carries forward is the vocabulary of the prior
template's prefatory framing and of the national platform. The Unification
Standard Level (USL) scores how ready a robot platform is; the Physical AI
Standard Level (PSL) scores how ready a trial site is; and the national MCP-PAI
standard supplies the audit, ledger, and provenance infrastructure that makes a
verification record trustworthy after the fact. The national platform
\cite{kawchak2026national} demonstrates these standards at the scale of a working
trial, and the author's prior survey of proposed United States bills
\cite{kawchak2026bills} establishes the patient-priority framing that this
instrument adopts: the patient, not the platform vendor and not the sponsor, is
the party the verification gate protects.

\subsection{Why a Law, and Why Now}

A reasonable reader will ask whether existing authority already reaches this
problem, and whether a new statute is premature for a field this young. The
national platform's executive summary answers both questions, and the bill adopts
its posture.\ On the first question, the existing frameworks can be adapted rather
than replaced. The Federal Food, Drug, and Cosmetic Act, the investigational
pathways, the Medicare medical-necessity standard, the privacy and
electronic-records rules, and the consensus engineering standards already supply
most of the substantive content a verification gate must enforce; what they do
not supply is the sequencing rule that makes verification a precondition of
generation rather than an afterthought. This bill adds that one rule and leaves
the adapted frameworks in place, a far less disruptive intervention than a new
regulatory regime would be.

On the second question, the timing is not premature but overdue. Physical AI
complexity has already passed the point at which manual, document-by-document
review can keep pace with the volume of machine-generated control code that a
multi-site trial produces, and the gap widens with every robot added to a site.
At the same time, the Food and Drug Administration's 2026 initiative to implement
real-time clinical trials \cite{fda2026realtime} makes faster, automated
assurance not merely tolerable but desirable, because a real-time trial cannot
wait weeks for a manual code review between procedures. 

The bill is written to
serve that initiative rather than to constrain it: its gate is automated, its
thresholds are bound to the agency's own credibility framework
\cite{fda2023credibility}, and its human-review and good-clinical-practice
requirements track the harmonised guidance the agency already recognises
\cite{iche6r3}. The purpose throughout is patient safety and efficacy, expressed
as a single operational rule. Mentions of the Food and Drug Administration and of
other bodies in this bill are respectful and opportunistic in the proper sense:
the United States can lead the world in surgical-humanoid VVUQ if it sets the
standard first, and this bill is offered to help it do so.

\subsection{Repository-Based AI as a Legislative Accelerator}

The second half of the bill's thesis concerns how this very document was made,
because the method is itself part of the argument. State-of-the-art
repository-based artificial intelligence models can read a corpus of technical
evidence, generate and execute verification code against it, and synthesise the
results into a public, reproducible document. This bill is such a document. Its
findings are not assertions; they are the recorded outputs of pipelines that any
reader can re-run from a stated seed, deposited under the repositories listed in
Table~\ref{tab:policy-sources} and citable by DOI \cite{repo-cancer-automated,
repo-physical-ai}.

The synthesis pattern is drawn from two places.\ The national platform's catalog
of source documents shows how adapted consensus standards and open repositories
can be turned into legislative blueprints rather than left as inert technical
artifacts. The author's second VVUQ study supplies the operational lineage that
this bill follows: an instruction is written, code is generated from
it, the code is executed under an assurance gate, and the run is synthesised into
a paper. 

The coding agent that performs those steps \cite{anthropicclaudecode}
makes the path from technical evidence to statutory text short enough that the
evidence and the law can be revised together. The practical consequence for the
Committee is that the record behind H.R. 9510 is not a static appendix but a
living, re-executable corpus, and the bill's documentation and attestation
sections are written so that any sponsor can produce the same kind of record for
its own robot-patient code.

\begin{table}[ht]
\centering
\begin{tabularx}{\textwidth}{L{7.5cm} Y}
\toprule
\textbf{Source file (abbreviated leaf)} & \textbf{What the section takes from it} \\
\midrule
\multicolumn{2}{l}{\textbf{papers/VVUQ-03/template-paper/}} \\
chunk\_01\_front\_matter.md & Prefatory-note tone; USL standard, PSL standard, MCP-PAI vocabulary \\
\multicolumn{2}{l}{\textbf{papers/VVUQ-01/final-paper/sections/}} \\
introduction.tex & Converging fronts; cost asymmetry; verification-first thesis \\
\multicolumn{2}{l}{\textbf{physical-ai .../final\_paper/sections/}} \\
executive\_summary.tex & Adapt-not-replace argument; economic and governance framing \\
source\_documents.tex & How adapted standards and repositories become legislative blueprints \\
\bottomrule
\end{tabularx}
\caption{Sources for the policy memorandum and concept paper.}
\label{tab:policy-sources}
\end{table}
